
Cross-architecture property prediction from neural network weights has been constrained by incompatible weight structures and architectural heterogeneity. Prior methods chose between architecture-specific encoding that preserves structure but requires same-family architectures, and architecture-invariant encoding that achieves cross-family transfer but discards operation-specific signals.

CAPE decomposes architectural differences into three learnable signals: operation-specific weight patterns, task-aligned embeddings, and computational graph topology. In proof-of-concept validation with 100 models and synthetic accuracy labels, CAPE achieved $\rho$ = 0.67 on $ResNet\rightarrow ViT$ transfer, a 24\% improvement over the SNE baseline ($\Delta\rho$ = 0.13, p = 0.032). All three components contributed independently. Diagnostic metrics confirmed distinct operation representations, preserved architectural structure, and meaningful topology contribution.

These results establish that architectural differences are not insurmountable barriers but separable signals amenable to modular learning. The proof-of-concept validates mechanism function; full-scale validation with 400 models across four architectures and real accuracy labels is required to demonstrate predictive capability for practical applications including automated model selection and architecture search.
